\section{Mechanism Synthesis and Open-Question Specifications}\label{sec:mechanism-agenda}
The main synthesis combines observations already recorded in the forty-four-setting matrix. A source result, a deduction about the protection scope, and a proposed mechanism test have different roles. Table~\ref{tab:mechanism-traceability} gives the links. The review package's \texttt{mechanism-synthesis/} files also connect each conditional-risk row and open question to setting IDs. Those links make the interpretation inspectable; they add no experimental observation.

\begin{table*}[t]\centering\small
\caption{Source relationships used in the synthesis. Execution changes qualify both feedback exposure and cover evaluation. Setting IDs refer to the source matrix and its version-specific passage locators.}
\label{tab:mechanism-traceability}
\setlength{\tabcolsep}{4pt}
\begin{tabularx}{\textwidth}{@{}P{.21\textwidth}P{.42\textwidth}Y@{}}\toprule
Finding & Observed basis & Conditional implication and open test \\\midrule
Execution organization & CG-composition and CB-topology report different transfer changes. TR-scaffold concerns terminal model attribution; ATR-main uses active repairs. TR-within and target/outcome records distinguish excluded sessions from observable step errors. & Recurring-task privacy and changed-deployment monitoring need distinct tests. Q1 asks which feature families survive specified changes; error normalization is an untested defense hypothesis. \\\addlinespace[3pt]
Distributed exposure & WRA-prompt, CB-task and ANT-model-scope concern different interfaces. PT-metadata/PT-content compare a restricted view with added arguments. Recipient and active settings require their stated rights. & Enumerate the representations available to the observer. Q2 isolates feedback; Q4 tests discovery and association. The local complete-model scores do not supply feedback attribution. \\\addlinespace[3pt]
Cover and execution cost & WRA-decoy and WRA-visibility have different utility conditions. CG-live pairs protected refitting with transport cost. The paired feedback study reports phase and complete-execution cells separately. & Constrain original-task utility when cover becomes work. Q3 tests transport termination versus isolated execution; harmful writes are not attributed to WRA. \\
\bottomrule\end{tabularx}\end{table*}

\subsection{Q1: feature stability under execution changes}
Hold the private label, observer rights and measurement layer fixed before comparing an execution change. Split by the changed composition or implementation, and keep repeated runs from the same task or session together. Frozen transfer and local adaptation answer different questions and should use separate results. Retain failures and refusals as measured outcomes, with any eligibility restrictions declared. A joint holdout tests interactions that separate single-axis tests cannot settle. Analysis should identify which feature family changes its predictive utility; an accuracy loss alone does not identify memorization as its cause.

\subsection{Q2: feedback attribution}
Use the same initial request for every category and disclose category-dependent information only in the later tool response. Retain prefix measurements as a negative control, predeclare the primary attacker/comparison and minimum effect of interest, and separate task failures from missing responses. Section~\ref{supp:randomized-feedback} implements this design with paired full and projected forwarding. The projection removes supplementary records while retaining the task fields; its effect combines changes in content, encoded length and subsequent execution. Separating these mediators requires additional interventions. Where a forced call path would invalidate the task, measure that path as a mediator rather than calling the comparison content-only. The artifact archive retains the earlier common-prefix population and its inconclusive result.

\subsection{Q3: cover termination and semantic isolation}
Define the original task's required final state, allowable additional work and resource budget before evaluating cover. A transport-terminated cover event and an executed decoy are different interventions. Isolation should be checked for business-state writes, quota consumption and feedback reaching the planner; a hidden response is not a rollback. Compare protection, overhead and original-task completion on matched attempts, including rejected or interrupted attempts with their actual status. Any state or utility tolerance and any equivalence criterion must be specified before interpreting a small difference as preservation.

\subsection{Q4: discovery, association and inference}
Record live concurrent execution and retain connection/stream association only when the chosen observer can obtain it. Keep complete-task and visible-tool intervals distinct. Use ground-truth clocks and event logs for evaluation, not to supply test-time task boundaries. Select detection thresholds on separate workloads or loads; report precision--recall and false alarms per time alongside correctly localized classification. Include a tool-only observer and an observer that can associate model and tool traffic. Section~\ref{supp:randomized-feedback} tests live concurrent execution with varying session composition, empty sessions and repeated classes on one host. Its IoU sensitivity and constant controls test whether a detector resolves instances or merely covers coarse activity. The earlier fixed-composition study is in the artifact archive. Transfer across workload, load and production background traffic remains open; the offline compositions are a separate detector diagnostic.
